% !PRINTED 431
\weberpage{431}
%<<CONT
setzt. Die Gleichung (6) hat dieselbe \gesperrt{Galois}'sche Gruppe, wie (7). Wir können nun eine beliebige unter den Wurzeln von (7) für $U_\infty$ setzen, und folglich können wir $U_\infty$ und $v_\infty$ in (9) einander entsprechen lassen. Dann können wir von den übrigen Wurzeln von (7) eine beliebige für $U_0$ nehmen, weil wir $x_1, x_2$ in \S.\,113, (4) durch $\varepsilon^r x_1$, $\varepsilon^{-r} x_2$ ersetzen können. Wir lassen also $U_0$ und $v_0$ einander entsprechen. Endlich können wir über die imaginäre fünfte Einheitswurzel $\varepsilon$ noch so verfügen, dass $U_1$ und $v_1$ sich entsprechen, und dann folgt aus der Uebereinstimmung der Gruppe, dass überhaupt $U_\xi$ und $v_\xi$ sich entsprechen, dass also
\weq{10}{U_\xi = -\gamma_2(\omega)v_\xi, \qquad \xi = \infty, 0, 1, 2, 3, 4}
ist. \gesperrt{Wenn wir also} $\omega$ \gesperrt{bei gegebenem} $z$ \gesperrt{aus der transcendenten Gleichung (8) bestimmen, so geben uns (10) und (5) die vollständige Lösung der Ikosaëderresolvente} 6\textsuperscript{ten} \gesperrt{Grades.}

Nach (2) und (5) ergiebt sich für $\xi = 0, 1, 2, 3, 4$
\weq{11}{\eta(\omega)\sqrt{v_\xi} = \sum_{-\infty,\infty}^{\nu}(-1)^\nu e^{\frac{\pi i}{60}(\omega - 24\xi)(6\nu+1)^2}.}
Der Exponent $(6\nu+1)^2$ kann nach dem Modul 5 nur den Rest $0, 1, -1$ geben. Lässt man also $\nu_1, \nu_2$ alle Zahlen durchlaufen, die den Bedingungen
\[
\nu_1 \equiv 1, 2, \qquad \nu_2 \equiv 0, 3 \pmod{5}
\]
genügen, so können wir den Ausdruck (11) so zerlegen:
\weq{12}{\eta(\omega)\sqrt{v_\xi} = \eta(5\omega) + \varepsilon^\xi \sum^{\nu_1}(-1)^{\nu_1} e^{\frac{\pi i}{60}\omega(6\nu_1+1)^2} + \varepsilon^{-\xi}\sum^{\nu_2}(-1)^{\nu_2} e^{\frac{\pi i}{60}\omega(6\nu_2+1)^2},}
wenn $\varepsilon = e^{\frac{2\pi i}{5}}$ gesetzt ist. Hieraus erkennt man, dass die Grössen $\sqrt{v_\xi}$ die Relationen (9) des vorigen Paragraphen befriedigen, und die Formel (10), \S.\,113 ergiebt für die Wurzel $x = x_1 : x_2$ der Ikosaëdergleichung den Ausdruck:
\weq{13}{\eta(5\omega)x = \sum^{\nu_1}(-1)^{\nu_1} e^{\frac{\pi i}{60}\omega(6\nu_1+1)^2}.}
Die übrigen Wurzeln erhält man, wenn man auf $x$ die Ikosaëdersubstitutionen anwendet. Man kann sie aber auch da-
%CONT>>
